% Options for packages loaded elsewhere
\PassOptionsToPackage{unicode}{hyperref}
\PassOptionsToPackage{hyphens}{url}
\documentclass[
  10pt,
  a4paper,
]{article}
\usepackage{xcolor}
\usepackage[margin=22mm]{geometry}
\usepackage{amsmath,amssymb}
\setcounter{secnumdepth}{-\maxdimen} % remove section numbering
\usepackage{iftex}
\ifPDFTeX
  \usepackage[T1]{fontenc}
  \usepackage[utf8]{inputenc}
  \usepackage{textcomp} % provide euro and other symbols
\else % if luatex or xetex
  \usepackage{unicode-math} % this also loads fontspec
  \defaultfontfeatures{Scale=MatchLowercase}
  \defaultfontfeatures[\rmfamily]{Ligatures=TeX,Scale=1}
\fi
\usepackage{lmodern}
\ifPDFTeX\else
  % xetex/luatex font selection
\fi
% Use upquote if available, for straight quotes in verbatim environments
\IfFileExists{upquote.sty}{\usepackage{upquote}}{}
\IfFileExists{microtype.sty}{% use microtype if available
  \usepackage[]{microtype}
  \UseMicrotypeSet[protrusion]{basicmath} % disable protrusion for tt fonts
}{}
\makeatletter
\@ifundefined{KOMAClassName}{% if non-KOMA class
  \IfFileExists{parskip.sty}{%
    \usepackage{parskip}
  }{% else
    \setlength{\parindent}{0pt}
    \setlength{\parskip}{6pt plus 2pt minus 1pt}}
}{% if KOMA class
  \KOMAoptions{parskip=half}}
\makeatother
\setlength{\emergencystretch}{3em} % prevent overfull lines
\providecommand{\tightlist}{%
  \setlength{\itemsep}{0pt}\setlength{\parskip}{0pt}}
\usepackage{mathrsfs}
\usepackage{mathtools}
\usepackage{xurl}
\emergencystretch=3em
\setcounter{tocdepth}{1}
\newcommand{\Beta}{\mathrm{B}}
\DeclareUnicodeCharacter{00B2}{\ensuremath{^2}}
\DeclareUnicodeCharacter{00B1}{\ensuremath{\pm}}
\DeclareUnicodeCharacter{00D7}{\ensuremath{\times}}
\DeclareUnicodeCharacter{2192}{\ensuremath{\to}}
\DeclareUnicodeCharacter{21A6}{\ensuremath{\mapsto}}
\DeclareUnicodeCharacter{2208}{\ensuremath{\in}}
\DeclareUnicodeCharacter{2209}{\ensuremath{\notin}}
\DeclareUnicodeCharacter{2212}{\ensuremath{-}}
\DeclareUnicodeCharacter{2227}{\ensuremath{\wedge}}
\DeclareUnicodeCharacter{2228}{\ensuremath{\vee}}
\DeclareUnicodeCharacter{2260}{\ensuremath{\ne}}
\DeclareUnicodeCharacter{2264}{\ensuremath{\le}}
\DeclareUnicodeCharacter{2265}{\ensuremath{\ge}}
\DeclareUnicodeCharacter{2297}{\ensuremath{\otimes}}
\DeclareUnicodeCharacter{00B3}{\ensuremath{^3}}
\DeclareUnicodeCharacter{2074}{\ensuremath{^4}}
\DeclareUnicodeCharacter{2076}{\ensuremath{^6}}
\DeclareUnicodeCharacter{03BB}{\ensuremath{\lambda}}
\DeclareUnicodeCharacter{03BC}{\ensuremath{\mu}}
\DeclareUnicodeCharacter{03B5}{\ensuremath{\epsilon}}
\DeclareUnicodeCharacter{03C4}{\ensuremath{\tau}}
\DeclareUnicodeCharacter{03A6}{\ensuremath{\Phi}}
\DeclareUnicodeCharacter{03B1}{\ensuremath{\alpha}}
\DeclareUnicodeCharacter{03C3}{\ensuremath{\sigma}}
\DeclareUnicodeCharacter{03B4}{\ensuremath{\delta}}
\DeclareUnicodeCharacter{03C0}{\ensuremath{\pi}}
\DeclareUnicodeCharacter{039B}{\ensuremath{\Lambda}}
\DeclareUnicodeCharacter{2202}{\ensuremath{\partial}}
\DeclareUnicodeCharacter{0393}{\ensuremath{\Gamma}}
\DeclareUnicodeCharacter{2205}{\ensuremath{\varnothing}}
\DeclareUnicodeCharacter{221E}{\ensuremath{\infty}}
\DeclareUnicodeCharacter{2200}{\ensuremath{\forall}}
\DeclareUnicodeCharacter{00B9}{\ensuremath{^1}}
\DeclareUnicodeCharacter{00B0}{\ensuremath{{}^\circ}}
\DeclareUnicodeCharacter{211D}{\ensuremath{\mathbb{R}}}
\usepackage{bookmark}
\IfFileExists{xurl.sty}{\usepackage{xurl}}{} % add URL line breaks if available
\urlstyle{same}
\hypersetup{
  hidelinks,
  pdfcreator={LaTeX via pandoc}}

\author{}
\date{}

\begin{document}

{
\setcounter{tocdepth}{3}
\tableofcontents
}
\section{The localized Weil packet: values, jets, and the observation
multiplier}\label{the-localized-weil-packet-values-jets-and-the-observation-multiplier}

This paper computes the finite Hermitian form specified by the zero side
of the Weil pairing on the programme's entire functions. It keeps the
polynomial coordinate change, all primary jets, and the actual
multiplier \(U=M_{j_h(v_h)}\). The form has a four-dimensional
nondegenerate quotient and a fully described nilpotent radical. Every
jet remains present in the source algebra and in the multiplier,
although the zero-side pairing evaluates only the zeroth jets.

The definitions of the quartet, \(E_h\), the Taylor remainder \(j_h\),
the unit \(U\), and the primary basis are those of the programme's OPG1
and OPG10--OPG11 in \emph{Full observation proofs}. They are restated
and proved to the extent needed here. The starting suggestion is
Proposition 3 in the supplied \emph{Split-Zero programme: orientation
and new angles}, 22 September 2026. The formula there changes from the
Fourier coordinate to the \(s\)-coordinate without changing the
polynomial's name. Equations (WP5)--(WP11) below give the exact
coordinate maps.

No analytic explicit-formula identity, Fourier inversion assertion, or
positivity theorem is assumed in the finite calculation. Section 10
specifies the precise zero-side identity proved here and its analytic
scope.

\subsection{1. The complete-order quartet and the two polynomial
coordinates}\label{the-complete-order-quartet-and-the-two-polynomial-coordinates}

Retain the programme's hypotheses \[
\rho=\frac12+\delta+i\gamma,
\qquad 0<\delta<\frac12,\quad \gamma>2,\quad m\in\mathbb Z_{\ge1}.
\tag{WP1}
\] Order its four distinct points as \[
\alpha_1=\rho,\quad
\alpha_2=1-\overline\rho,\quad
\alpha_3=\overline\rho,\quad
\alpha_4=1-\rho,
\qquad \mathcal R=\{\alpha_1,\alpha_2,\alpha_3,\alpha_4\}.
\tag{WP2}
\] The real parts distinguish points with the same imaginary part
because \(\delta>0\), and the imaginary parts distinguish the upper and
lower pairs because \(\gamma>0\). Let \(g=2\xi\) have a zero of exact
order \(m\) at each \(\alpha\in\mathcal R\). Put \[
d(s)=\prod_{\alpha\in\mathcal R}(s-\alpha),\qquad
h(s)=d(s)^m,\qquad n=4m,\qquad v_h(s)=\frac{g(s)}{h(s)}.
\tag{WP3}
\] Here \(s\) is the original zeta coordinate. At each \(\alpha\), write
\(g(s)=(s-\alpha)^m b_\alpha(s)\), where \(b_\alpha\) is holomorphic and
\(b_\alpha(\alpha)\ne0\). If \(h_\alpha(s)=h(s)/(s-\alpha)^m\), then
\(h_\alpha(\alpha)\ne0\), and \[
v_h(s)=\frac{b_\alpha(s)}{h_\alpha(s)}\quad\hbox{near }\alpha,
\qquad
v_h(\alpha)=\frac{g^{(m)}(\alpha)}{m!\,h_\alpha(\alpha)}\ne0.
\tag{WP4}
\] Thus all the apparent singularities of \(g/h\) are removable and
\(v_h\) is entire. At every zero \(\rho'\notin\mathcal R\) of \(g\), it
vanishes because \(h(\rho')\ne0\).

Define the Fourier coordinate and its inverse by \[
\theta(s)=\frac{s-1/2}{i},\qquad \iota(z)=\frac12+iz,
\qquad \theta\iota=\mathrm{id}_{\mathbb C},\quad
\iota\theta=\mathrm{id}_{\mathbb C}.
\tag{WP5}
\] For \(P\in\mathbb C[z]\), the associated \(s\)-polynomial is \[
q_P(s)=P(\theta(s)),
\qquad F_P(z)=P(z)v_h(\iota(z)).
\tag{WP6}
\] In particular \(F_P(\theta(s))=q_P(s)v_h(s)\). The two polynomials
\(P\) and \(q_P\) have the same degree, but their coefficients and
variables differ by (WP5). A factor \(P(1-\overline\alpha)\) is
therefore not the value obtained from \(F_P\) unless \(P\) has first
been replaced by \(q_P\).

Put \(\beta_\alpha=\theta(\alpha)\). In the order (WP2), \[
(\beta_{\alpha_1},\beta_{\alpha_2},\beta_{\alpha_3},\beta_{\alpha_4})
=(\gamma-i\delta,\gamma+i\delta,-\gamma-i\delta,-\gamma+i\delta).
\tag{WP7}
\] Let \(\sigma\alpha=1-\overline\alpha\). This involution interchanges
\(\alpha_1,\alpha_2\) and \(\alpha_3,\alpha_4\). Direct substitution
gives \[
\overline{\theta(\alpha)}=\theta(\sigma\alpha),
\qquad \iota(\overline{\beta_\alpha})=\sigma\alpha.
\tag{WP8}
\]

Define \[
H(z)=\prod_{\alpha\in\mathcal R}(z-\beta_\alpha)^m
=i^{-n}h(\iota(z)),
\qquad E_z=\mathbb C[z]/(H),\quad E_h=\mathbb C[s]/(h).
\tag{WP9}
\] Here \(n=4m\), so \(i^{-n}=1\); the scalar is displayed to record how
the defining polynomials transform. The maps \[
\mathcal C:E_z\longrightarrow E_h,
\quad [P]\longmapsto[P\circ\theta],
\qquad
\mathcal C^{-1}:E_h\longrightarrow E_z,
\quad[q]\longmapsto[q\circ\iota]
\tag{WP10}
\] are inverse algebra isomorphisms. Indeed substitution sends \(H\) to
\(i^{-n}h\) and \(h\) to \(i^nH\), so the maps descend to the quotients.
Their two composites fix every polynomial by (WP5). They also retain the
multiplication operators exactly: \[
\mathcal C M_z=M_{\theta(s)}\mathcal C,
\qquad
\mathcal C M_{\iota(z)}=M_s\mathcal C.
\tag{WP11}
\]

\subsection{2. All the jets and the entire-function
remainder}\label{all-the-jets-and-the-entire-function-remainder}

For each \(\alpha\in\mathcal R\), introduce an independent nilpotent
coordinate \(t_\alpha\) and the algebra \[
A_\alpha=\mathbb C[t_\alpha]/(t_\alpha^m).
\tag{WP12}
\] There is an algebra isomorphism \[
\mathcal J:E_h\longrightarrow\prod_{\alpha\in\mathcal R}A_\alpha,
\qquad
[q]\longmapsto\left(
\sum_{a=0}^{m-1}\frac{q^{(a)}(\alpha)}{a!}t_\alpha^a
\right)_{\alpha\in\mathcal R}.
\tag{WP13}
\] Every derivative in this formula is part of the specified map.

Here is a direct proof including its inverse. Define \[
T_\alpha(s)=\sum_{a=0}^{m-1}
\frac{(1/h_\alpha)^{(a)}(\alpha)}{a!}(s-\alpha)^a,
\qquad e_\alpha=[h_\alpha T_\alpha]\in E_h.
\tag{WP14}
\] Modulo \((s-\alpha)^m\), the product \(h_\alpha T_\alpha\) equals
\(1\), by multiplication of the Taylor expansion and its truncated
inverse. Modulo \((s-\eta)^m\) for \(\eta\ne\alpha\), it equals \(0\),
because \(h_\alpha\) is divisible by that factor. Thus
\(\mathcal J(e_\alpha)\) is \(1\) in the \(\alpha\)-factor and \(0\)
elsewhere. The inverse of (WP13) is \[
\mathcal J^{-1}\left(\left(\sum_{a=0}^{m-1}c_{\alpha,a}t_\alpha^a\right)_\alpha\right)
=\sum_{\alpha\in\mathcal R}\sum_{a=0}^{m-1}
c_{\alpha,a}e_\alpha(s-\alpha)^a.
\tag{WP15}
\] This proves surjectivity. For injectivity, a polynomial in the kernel
has a zero of order at least \(m\) at each of the four distinct points.
Successive division by their linear factors shows it is divisible by
their product \(h\). The kernel is therefore zero in \(E_h\). Taylor
multiplication modulo \(t_\alpha^m\) proves multiplicativity. Equations
(WP13)--(WP15) also prove \[
e_\alpha e_\eta=0\ (\alpha\ne\eta),\qquad
e_\alpha^2=e_\alpha,\qquad\sum_\alpha e_\alpha=1,
\tag{WP16}
\] and give the basis \(e_\alpha(s-\alpha)^a\), \(0\le a<m\), of
\(E_h\).

For an entire function \(f\), define its finite Taylor remainder by \[
j_h(f)=\mathcal J^{-1}\left(\left(
\sum_{a=0}^{m-1}\frac{f^{(a)}(\alpha)}{a!}t_\alpha^a
\right)_\alpha\right)\in E_h.
\tag{WP17}
\] This is an algebra homomorphism: addition and scalar multiplication
follow coefficient by coefficient, and the Leibniz rule gives
multiplication of the truncated Taylor series. It agrees with the
quotient map on polynomials. The class \(j_h(f)\) has a unique
representative of degree less than \(n\), by division by the monic
polynomial \(h\).

In the programme notation \(j_hv_h\), the symbol \(j_h\) denotes this
map applied to \(v_h\). It is not a second scalar function multiplied by
\(v_h\). Put \[
u_h=j_h(v_h),\qquad U=M_{u_h}:E_h\longrightarrow E_h.
\tag{WP18}
\] In the \(\alpha\)-factor of (WP13), set \[
v_{\alpha,a}=\frac{v_h^{(a)}(\alpha)}{a!},
\qquad u_\alpha(t_\alpha)=\sum_{a=0}^{m-1}v_{\alpha,a}t_\alpha^a.
\tag{WP19}
\] All constants \(v_{\alpha,0}\) are nonzero by (WP4). Write
\(u_\alpha=v_{\alpha,0}(1+b_\alpha)\), with \(b_\alpha\in(t_\alpha)\).
Since \(b_\alpha^m=0\), its inverse is \[
u_\alpha^{-1}=v_{\alpha,0}^{-1}
\sum_{r=0}^{m-1}(-b_\alpha)^r.
\tag{WP20}
\] Multiplication of this finite series by \(u_\alpha\) gives
\(1-(-b_\alpha)^m=1\). Hence \(u_h\) is a unit and \(U\) an
automorphism. In full jet coordinates its formula is \[
(Ux)_{\alpha,a}=\sum_{b=0}^a v_{\alpha,b}x_{\alpha,a-b}
\qquad(0\le a<m).
\tag{WP21}
\] It therefore keeps every derivative of the amplitude through order
\(m-1\). The determinant in the primary basis is
\(\prod_{\alpha\in\mathcal R}v_h(\alpha)^m\ne0\), because each block in
(WP21) is triangular with that constant diagonal.

\subsection{3. The zeroth-jet quotient and its exact multiplier
map}\label{the-zeroth-jet-quotient-and-its-exact-multiplier-map}

Let \(V_{\mathcal R}=\mathbb C^{\mathcal R}\), with coordinatewise
multiplication, and define \[
e:E_h\longrightarrow V_{\mathcal R},\qquad
e([q])=(q(\alpha))_{\alpha\in\mathcal R}.
\tag{WP22}
\] This is a surjective algebra homomorphism: it is the constant-term
map in each factor of (WP13), and (WP15) lifts any tuple by
\(\sum c_\alpha e_\alpha\). Its kernel is \[
J=\ker e=(d)/(d^m)\subset E_h,
\qquad\dim_{\mathbb C}J=4m-4.
\tag{WP23}
\] Indeed a polynomial has value zero at every \(\alpha\) precisely when
it is divisible by \(d\). In the primary coordinates, \(J\) is
\(\prod_\alpha(t_\alpha)\); it has basis \(e_\alpha(s-\alpha)^a\),
\(1\le a<m\), and the asserted dimension follows. Every element of \(J\)
has its \(m\)-th power zero. Conversely, a nilpotent element must have
value zero in each complex coordinate, since a complex number with a
vanishing positive power is zero. Thus \(J\) is exactly the nilradical
of \(E_h\).

Write \[
V_h:V_{\mathcal R}\longrightarrow V_{\mathcal R},\qquad
(c_\alpha)_\alpha\longmapsto(v_h(\alpha)c_\alpha)_\alpha,
\quad
a_h=eU:E_h\longrightarrow V_{\mathcal R}.
\tag{WP24}
\] Then \[
eU=V_he,\qquad \ker a_h=J,
\qquad a_h\text{ is surjective}.
\tag{WP25}
\] The identity follows either from evaluation of a product or from the
\(a=0\) case of (WP21). The map \(V_h\) is invertible because its four
diagonal entries are nonzero, so the kernel and surjectivity assertions
follow from (WP22)--(WP23). Thus the exact comparison with the
programme's multiplier is the commuting diagram of typed maps \[
\begin{array}{ccc}
E_h&\xrightarrow{\ U\ }&E_h\\
e\downarrow&&\downarrow e\\
V_{\mathcal R}&\xrightarrow{\ V_h\ }&V_{\mathcal R}.
\end{array}
\tag{WP26}
\] The left and right vertical maps retain the same kernel \(J\); the
horizontal map on the upper row retains all the additional coefficients
(WP21).

The entire-function family is a further, distinct vector space \[
\mathcal F_h=\{F_P:P\in\mathbb C[z]\}.
\tag{WP26a}
\] The map \(P\mapsto F_P\) is a linear isomorphism onto this family.
Surjectivity is its definition. For injectivity, \(v_h\) is nonzero at a
packet root, hence nonzero on a neighborhood of that root. If \(F_P\)
vanishes identically, \(P\) vanishes on the corresponding open set in
the \(z\)-coordinate. A nonzero polynomial has only finitely many roots,
so \(P=0\).

The full-jet map from this entire-function family is \[
\begin{aligned}
\mathcal T_h:\mathcal F_h&\longrightarrow E_h,
&\mathcal T_h(F)&=j_h(F\circ\theta),\\
\mathcal T_h(F_P)&=U\mathcal C[P],
&(\mathcal J\mathcal T_h(F))_{\alpha,a}
&=i^{-a}\frac{F^{(a)}(\beta_\alpha)}{a!}.
\end{aligned}
\tag{WP26b}
\] The first identity on the second line follows from (WP6), (WP17), and
(WP18). The second follows by differentiating the affine change of
variable \(\theta\) exactly \(a\) times; its derivative is \(1/i\). Thus
the factor \(i^{-a}\) is retained in every jet coordinate.

The map \(\mathcal T_h\) is onto and its kernel is \[
\ker\mathcal T_h=\{F_{HQ}:Q\in\mathbb C[z]\}.
\tag{WP26c}
\] Indeed \(U\) and \(\mathcal C\) are isomorphisms, so
\(\mathcal T_h(F_P)=0\) exactly when \(P\) is divisible by \(H\). There
is a separate surjection \[
e\mathcal T_h:\mathcal F_h\longrightarrow V_{\mathcal R},
\quad F\longmapsto(F(\beta_\alpha))_\alpha,
\qquad
\ker(e\mathcal T_h)=\{F_{D_zQ}:Q\in\mathbb C[z]\},
\quad D_z(z)=\prod_\alpha(z-\beta_\alpha).
\tag{WP26d}
\] The coordinate formula follows from (WP26b) at \(a=0\). It is onto by
(WP25). Since \(v_h(\alpha)\ne0\), its value at \(F_P\) is zero exactly
when \(P\) vanishes at all the distinct \(\beta_\alpha\), which is
equivalent to divisibility by \(D_z\). Thus the entire functions
themselves are not identified with their finite jets or with their
values: (WP26b)--(WP26d) give the actual surjections and both exact
kernels.

\subsection{4. The finite Hermitian form and its
radical}\label{the-finite-hermitian-form-and-its-radical}

Use the convention that a Hermitian form is conjugate-linear in its
first argument and linear in its second. On \(V_{\mathcal R}\), define
\[
b_{\mathcal R}(c,c')=m\sum_{\alpha\in\mathcal R}
\overline{c_{\sigma\alpha}}c'_\alpha.
\tag{WP27}
\] It is Hermitian: the conjugate of \(b_{\mathcal R}(c',c)\), after
replacing the summation index \(\alpha\) by \(\sigma\alpha\), is exactly
(WP27). It is nondegenerate. If it vanishes against every \(c'\),
selecting \(c'\) with a single nonzero coordinate at \(\alpha\) gives
\(m\overline{c_{\sigma\alpha}}=0\), hence every coordinate of \(c\) is
zero.

Define the packet form on the full algebra by the pullback \[
B_h(x,y)=b_{\mathcal R}(a_hx,a_hy)
=m\sum_{\alpha\in\mathcal R}
\overline{v_h(\sigma\alpha)x(\sigma\alpha)}
v_h(\alpha)y(\alpha),
\qquad x,y\in E_h.
\tag{WP28}
\] Evaluation is independent of polynomial representatives by (WP22), so
the formula is well defined. It is Hermitian because (WP27) is. Let
\(\operatorname{rad}B_h\) mean the set of \(x\) with \(B_h(x,y)=0\) for
every \(y\in E_h\).

\textbf{Theorem WP1 (exact radical and quotient).} There is an exact
sequence \[
0\longrightarrow J\longrightarrow E_h
\xrightarrow{\ a_h\ }V_{\mathcal R}\longrightarrow0,
\qquad \operatorname{rad}B_h=J.
\tag{WP29}
\] The induced map \(\overline a_h:E_h/J\to V_{\mathcal R}\) is an
algebra-multiplier comparison and a linear isometry from the descended
Hermitian form to \(b_{\mathcal R}\). More precisely,
\(e:E_h/J\to V_{\mathcal R}\) is an algebra isomorphism, and
\(\overline a_h=V_h\overline e\) is a linear isomorphism;
\(\overline a_h\) need not preserve the unit or multiplication.

\textbf{Proof.} Exactness is (WP25). If \(x\in J\), then \(a_hx=0\), so
(WP28) vanishes for every \(y\). Conversely, if (WP28) vanishes for
every \(y\), surjectivity of \(a_h\) implies that
\(b_{\mathcal R}(a_hx,c')=0\) for every \(c'\). Nondegeneracy of (WP27)
gives \(a_hx=0\), hence \(x\in J\). The quotient map is an isomorphism
by exactness. Formula (WP28) proves the isometry identity. The
distinction between \(\overline e\) and \(\overline a_h\) follows from
(WP24): multiplying all coordinates by a fixed nonconstant unit is a
linear automorphism, and is not in general an algebra homomorphism.
\(\square\)

This theorem does not remove the nilpotent directions from the original
packet. It identifies the precise quotient seen by this particular
Hermitian form, and identifies every omitted direction as an element of
the specified ideal \(J\).

\subsection{5. Complete inertia and explicit negative
representatives}\label{complete-inertia-and-explicit-negative-representatives}

In the order (WP2), write \(c=(a,b,c_3,c_4)\). Then \[
b_{\mathcal R}(c,c')
=m\bigl(\overline b a'+\overline a b'
+\overline{c_4}c_3'+\overline{c_3}c_4'\bigr).
\tag{WP30}
\] The matrix is the direct sum of two blocks
\(m\begin{pmatrix}0&1\\1&0\end{pmatrix}\). In each pair the two vectors
\((1,1)/\sqrt{2m}\) and \((1,-1)/\sqrt{2m}\) have squared values \(+1\)
and \(-1\), and their mutual pairing is zero. They form a basis of that
two-dimensional pair. Consequently \[
\operatorname{inertia}(b_{\mathcal R})=(2,2,0),
\qquad
\operatorname{inertia}(B_h)=(2,2,4m-4),
\tag{WP31}
\] where inertia lists positive, negative, and radical dimensions in
that order. This is inertia under congruence of Hermitian forms, rather
than an assertion that an arbitrary coefficient-frame matrix of \(B_h\)
has eigenvalues \(\pm m\).

For a complete diagonal basis of the original packet, use \[
x_1^{\pm}=\frac1{\sqrt{2m}}U^{-1}(e_{\alpha_1}\pm e_{\alpha_2}),
\qquad
x_2^{\pm}=\frac1{\sqrt{2m}}U^{-1}(e_{\alpha_3}\pm e_{\alpha_4}),
\tag{WP32}
\] together with \[
x_{\alpha,a}=U^{-1}\bigl(e_\alpha(s-\alpha)^a\bigr),
\quad \alpha\in\mathcal R,\quad 1\le a<m.
\tag{WP33}
\] These form a basis because the original primary vectors do, the two
changes of basis at order zero are invertible, and \(U\) is invertible.
The last \(4m-4\) vectors lie in \(J\). Formula (WP28) gives
\(B_h(x_j^+,x_j^+)=1\), \(B_h(x_j^-,x_j^-)=-1\), all pairings between
distinct displayed positive or negative vectors zero, and all pairings
with (WP33) zero. This proves (WP31) on the entire vector space,
including all primary multiplicities.

These representatives are explicit finite Taylor expressions. For
example \[
U^{-1}e_\alpha
=e_\alpha\sum_{a=0}^{m-1}
\frac{(1/v_h)^{(a)}(\alpha)}{a!}(s-\alpha)^a.
\tag{WP34}
\] The reciprocal is taken only in a neighborhood of \(\alpha\), where
it is holomorphic by (WP4). Taylor multiplication proves (WP34), as also
follows from (WP20). Formula (WP34) prescribes the full jet of a
representative whose amplitude has constant primary jet \(1\) at one
root and zero at the others.

There are also degree-three representatives, without imposing any
condition on their higher amplitude jets. Define the Lagrange
polynomials \[
\ell_\alpha(s)=\prod_{\eta\in\mathcal R\setminus\{\alpha\}}
\frac{s-\eta}{\alpha-\eta}.
\tag{WP35}
\] Their denominators are nonzero, and
\(\ell_\alpha(\eta)=\delta_{\alpha\eta}\). Put \[
q_-(s)=\frac{\ell_{\alpha_1}(s)}{v_h(\alpha_1)}
-\frac{\ell_{\alpha_2}(s)}{v_h(\alpha_2)},
\qquad P_-(z)=q_-(1/2+iz).
\tag{WP36}
\] Then \(a_h[q_-]=(1,-1,0,0)\), so \[
B_h([q_-],[q_-])=-2m.
\tag{WP37}
\] Dividing either polynomial in (WP36) by \(\sqrt{2m}\) gives value
\(-1\); dividing by \(\sqrt m\) gives value \(-2\). The second negative
direction is obtained with \(\alpha_3,\alpha_4\). This formula includes
the exact Fourier-coordinate representative \(P_-\).

For \(m>1\), the class \([q_-]\) in (WP36) and the full-jet
representative \(U^{-1}(e_{\alpha_1}-e_{\alpha_2})\) need not agree.
Their difference lies in \(J\), because their images under \(a_h\) agree
and \(\ker a_h=J\). Thus the two explicit choices have the same pairings
with every packet class, with their difference retained as a specified
nilpotent class.

\subsection{6. The surviving jet filtration and multiplication
operators}\label{the-surviving-jet-filtration-and-multiplication-operators}

The radical has the finite filtration \[
E_h=J^0\supset J^1\supset\cdots\supset J^{m-1}\supset J^m=0,
\qquad J^r=(d^r)/(d^m).
\tag{WP38}
\] In primary coordinates, \(J^r=\prod_\alpha(t_\alpha^r)\). To verify
this from \(d\), write \(d(s)=(s-\alpha)c_\alpha(s)\), with
\(c_\alpha(\alpha)\ne0\). The factor \(c_\alpha\) is a unit in
\(A_\alpha\), so the ideal generated by \(d^r\) in that primary factor
is exactly \((t_\alpha^r)\). The isomorphism (WP13) then proves (WP38)
for every \(r\).

For \(0\le r<m\), the exact layer map is \[
\lambda_r:J^r/J^{r+1}\longrightarrow V_{\mathcal R},
\qquad [x]\longmapsto
\left(\frac{x^{(r)}(\alpha)}{r!}\right)_\alpha.
\tag{WP39}
\] For \(r=0\), this means the map \(e\) on \(E_h/J\). A polynomial
representative of an element of \(J^r\) has all derivatives of orders
less than \(r\) zero at each root. Changing it by a multiple of \(h\)
changes none of the derivatives in (WP39), since \(r<m\). Changing its
class by \(J^{r+1}\) also leaves those coefficients unchanged. Formula
(WP13) shows that the kernel is exactly \(J^{r+1}\), and that the
elements \(e_\alpha(s-\alpha)^r\) map to the coordinate basis. Hence
(WP39) is an isomorphism and every layer has dimension \(4\).

All these layers are preserved by \(U\), its inverse, and \(A_h=M_s\).
Their induced operators are \[
\lambda_r\,\operatorname{gr}_r(U)=V_h\lambda_r,
\qquad
\lambda_r\,\operatorname{gr}_r(A_h)
=\operatorname{diag}(\alpha)\lambda_r.
\tag{WP40}
\] For the first equality, use (WP21): the coefficient of order \(r\) in
a product with an element whose lower coefficients are zero is
\(v_{\alpha,0}x_{\alpha,r}\). For the second, multiply its Taylor series
by \(\alpha+t_\alpha\); the extra term only enters the next layer.
Before passage to a layer, the complete formulas remain (WP21) and \[
A_h\bigl(e_\alpha(s-\alpha)^a\bigr)
=\alpha e_\alpha(s-\alpha)^a+e_\alpha(s-\alpha)^{a+1},
\tag{WP41}
\] where the last term is zero for \(a=m-1\). Thus the original Jordan
chains and the amplitude derivatives survive in the source and its
filtered maps. The form (WP28) reads only the layer \(r=0\); this
follows from its evaluated formula and does not assert that higher
derivatives are unavailable to the observation map.

\subsection{7. The exact adjoint relation for every polynomial
multiplier}\label{the-exact-adjoint-relation-for-every-polynomial-multiplier}

For \(q\in\mathbb C[s]\), define \[
q^\#(s)=\overline{q(1-\overline s)}.
\tag{WP42}
\] This is a polynomial: if \(q(s)=\sum_k c_ks^k\), then
\(q^\#(s)=\sum_k\overline{c_k}(1-s)^k\). It is an antilinear,
multiplicative involution. Since \(\sigma\) permutes \(\mathcal R\), \[
h^\#(s)=\prod_{\alpha\in\mathcal R}(1-s-\overline\alpha)^m
=(-1)^{4m}\prod_{\alpha\in\mathcal R}(s-\sigma\alpha)^m=h(s).
\tag{WP43}
\] It therefore descends to an antilinear algebra involution on \(E_h\).
Its action on all the jets is \[
(x^\#)_{\alpha,a}=(-1)^a\overline{x_{\sigma\alpha,a}}.
\tag{WP44}
\] Indeed, substitute \(s=\alpha+t\) into (WP42), expand \(q\) at
\(\sigma\alpha\), and conjugate its expansion at the increment
\(-\overline t\). This gives the displayed coefficient, including its
sign.

For the original completed zeta function, its functional identities are
\(g(1-s)=g(s)\) and \(g(\overline s)=\overline{g(s)}\). They imply
\(g^\#=g\). The quartet polynomial has exactly the same two identities:
conjugation permutes its four factors, and replacing \(s\) by \(1-s\)
gives the sign \((-1)^{4m}=1\) and the permutation
\(\alpha\mapsto1-\alpha\). Therefore \[
v_h^\#=v_h,\qquad v_h(\sigma\alpha)=\overline{v_h(\alpha)},
\qquad (j_hv_h)^\#=j_hv_h.
\tag{WP44a}
\] The first identity follows by division off the finite zero set of
\(h\) and then by continuity at the removable points (WP4). The second
is its evaluation, and the third follows from the signed jet formula
(WP44). No scalar factor accompanies these identities for the monic
degree-\(4m\) polynomial in (WP3).

\textbf{Theorem WP2 (multiplier adjoints).} For every \(q,x,y\in E_h\),
\[
B_h(M_qx,y)=B_h(x,M_{q^\#}y).
\tag{WP45}
\] In particular, if \(Y_h=(A_h-\frac12I)/i=M_{\theta(s)}\), then \[
B_h(A_hx,y)=B_h(x,(I-A_h)y),
\qquad B_h(Y_hx,y)=B_h(x,Y_hy).
\tag{WP46}
\]

\textbf{Proof.} At each root \(\alpha\),
\(q^\#(\alpha)=\overline{q(\sigma\alpha)}\). Inserting a factor
\(q(\sigma\alpha)\) in the conjugated first entry of (WP28) therefore
has exactly the same effect as inserting \(q^\#(\alpha)\) in its second
entry. This proves (WP45) term by term. For the coordinate polynomial,
\(s^\#=1-s\). For \(\theta(s)=-i(s-1/2)\), antilinearity and this
equality give \(\theta^\#(s)=i((1-s)-1/2)=-i(s-1/2)=\theta(s)\). These
substitutions prove both assertions in (WP46). \(\square\)

The operator \(Y_h\) is consequently self-adjoint with respect to the
specified Hermitian form. Its induced eigenvalues on \(E_h/J\) are
exactly (WP7), since (WP40) gives the diagonal multiplier
\((\alpha-1/2)/i\). The nonreal values
\(\gamma\pm i\delta,-\gamma\pm i\delta\) are paired by the two
hyperbolic blocks (WP30). There is no claim that self-adjointness for
this indefinite form makes those values real. Equations (WP40), (WP44),
and (WP46) give the exact relation between the full primary operator and
its value quotient.

In particular the original amplitude unit is itself self-adjoint for
this form, by (WP44a) and (WP45). For polynomial inputs, the entire
function \[
\mathcal A_{P,Q}(s)=q_P^\#(s)q_Q(s)v_h(s)^2
\tag{WP46a}
\] has the packet values \[
\mathcal A_{P,Q}(\alpha)
=\overline{q_P(\sigma\alpha)v_h(\sigma\alpha)}
q_Q(\alpha)v_h(\alpha).
\tag{WP46b}
\] This follows from \(q_P^\#(\alpha)=\overline{q_P(\sigma\alpha)}\) and
(WP44a). On the complement of the zeros of \(g\), the exact
logarithmic-derivative identity is \[
\frac{\xi'(s)}{\xi(s)}\mathcal A_{P,Q}(s)
=q_P^\#(s)q_Q(s)\frac{g'(s)g(s)}{h(s)^2}.
\tag{WP46c}
\] Here \(g=2\xi\), so \(\xi'/\xi=g'/g\); multiplication by \((g/h)^2\)
proves the formula with no additional factor of two. At a zero of \(g\)
outside the packet, the right side is holomorphic because \(h\) does not
vanish there. At \(\alpha\in\mathcal R\), writing
\(g=(s-\alpha)^mb_\alpha\) and \(h=(s-\alpha)^mh_\alpha\) shows that
\(g'g/h^2\) has the Laurent term \(m v_h(\alpha)^2/(s-\alpha)\). Thus
its residue after multiplication by \(q_P^\#q_Q\) is exactly
\(m\mathcal A_{P,Q}(\alpha)\). These are local meromorphic identities;
no contour estimate or prime-side identity is used to prove them.

\subsection{8. The original period observation and what each map
reads}\label{the-original-period-observation-and-what-each-map-reads}

The map \(a_h=eU\) applies the same full-jet unit as the original
observation and then takes values. To make the relation to the period
stage explicit, let \(\Pi:E_h\to\mathbb C^n\) be the programme's OPG4
linear map, defined on the unique degree-less-than-\(n\) representative
by its prescribed convergent period integrals. Retain its actual
contours, exponential weight, and basis; none is changed here. Put \[
\mathbf b_{\alpha,a}=\Pi\bigl(e_\alpha(s-\alpha)^a\bigr).
\tag{WP47}
\] For every \(x\in E_h\), linearity and (WP21) give the identity of
vectors in \(\mathbb C^n\) \[
\Pi Ux=
\sum_{\alpha\in\mathcal R}\sum_{a=0}^{m-1}
\mathbf b_{\alpha,a}
\sum_{b=0}^a v_{\alpha,b}x_{\alpha,a-b}.
\tag{WP48}
\] This proof uses a finite sum of the original primary classes; it does
not assign a period integral to an arbitrary entire function. The value
map, with the same coefficients before projection, is \[
a_hx=(v_{\alpha,0}x_{\alpha,0})_\alpha.
\tag{WP49}
\] Thus (WP48) and (WP49) are two exact maps out of the same source
after the same multiplication. In particular, for \(x=j_h(f)\), the
ring-homomorphism property (WP17) gives \[
Uj_h(f)=j_h(v_hf),\qquad
\Pi Uj_h(f)=\Pi j_h(v_hf),\qquad
a_hj_h(f)=(v_h(\alpha)f(\alpha))_\alpha.
\tag{WP50}
\] The middle formula evaluates \(\Pi\) on a finite Taylor remainder,
exactly as in OPG11. It is not the integral of the unreduced entire
function unless a separate equality proves that assertion.

There is also a precise test for whether any proposed observation
\(T:E_h\to W\), to a complex vector space \(W\), factors through the
value map \(a_h\). The factorization exists exactly when \(T(J)=0\);
when it exists its only possible formula is \[
\overline T(c)=T\left(U^{-1}\sum_{\alpha\in\mathcal R}c_\alpha e_\alpha\right).
\tag{WP51}
\] To prove this, the expression in parentheses has \(a_h\)-image \(c\)
by (WP15) and (WP24). If \(T(J)=0\), any two lifts with that image
differ by (WP29) in \(J\), so their \(T\)-images agree, and (WP51) gives
a well-defined linear factorization. Surjectivity of \(a_h\) forces
uniqueness. Conversely a factorization through \(a_h\) kills its kernel
\(J\). This computes the exact space of factorizing maps as the
injective image \[
\operatorname{Hom}_{\mathbb C}(V_{\mathcal R},W)
\xrightarrow{\ a_h^*\ }
\operatorname{Hom}_{\mathbb C}(E_h,W),
\quad \overline T\longmapsto\overline T a_h,
\quad
\operatorname{im}(a_h^*)=\{T:T(J)=0\}.
\tag{WP52}
\] It does not presume that the period map, or the later cyclic tensor
observation, lies in that subspace.

\subsection{9. Why the off-line and on-line cases have different packet
algebras}\label{why-the-off-line-and-on-line-cases-have-different-packet-algebras}

The hypotheses (WP1) guarantee four distinct roots. Substituting
\(\delta=0\) into a formula which assumes four distinct roots merges
\(\alpha_1\) with \(\alpha_2\) and \(\alpha_3\) with \(\alpha_4\). If
one retains the product of four factors literally, their orders double.
Exact zero order \(m\) would then no longer justify dividing \(g\) by
that new polynomial. The on-line packet must instead be defined from its
distinct roots and their actual orders.

For completeness, the general finite value calculation is as follows.
Let \(\mathcal A\) be a finite set of distinct complex points stable
under \(\sigma\alpha=1-\overline\alpha\). Assign positive integer orders
\(m_\alpha\), with \(m_{\sigma\alpha}=m_\alpha\), and suppose an entire
\(g\) has these exact orders there. Set \[
h_{\mathcal A}(s)=\prod_{\alpha\in\mathcal A}(s-\alpha)^{m_\alpha},
\qquad v_{\mathcal A}=g/h_{\mathcal A}.
\tag{WP53}
\] With the corresponding jet algebra, remainder multiplier, and
amplitude-value map, define \[
B_{\mathcal A}(x,y)=
\sum_{\alpha\in\mathcal A}m_\alpha
\overline{v_{\mathcal A}(\sigma\alpha)x(\sigma\alpha)}
v_{\mathcal A}(\alpha)y(\alpha).
\tag{WP54}
\] Let \(f\) be the number of fixed points of \(\sigma\) in
\(\mathcal A\), and let \(p\) be the number of its two-element orbits.
Then \[
\operatorname{inertia}(B_{\mathcal A})
=\left(f+p,\ p,\ \sum_{\alpha\in\mathcal A}(m_\alpha-1)\right).
\tag{WP55}
\] To prove this directly, use the local rings
\(\mathbb C[t_\alpha]/t_\alpha^{m_\alpha}\). The interpolation proof
(WP14)--(WP15) applies with the individual order in each factor, because
the other factors are still units at that root. Evaluation has kernel
\(\prod_\alpha(t_\alpha)\) of the stated dimension, and amplitude
multiplication is invertible because its constants are nonzero. On its
value quotient, a fixed point contributes the positive form
\(m_\alpha|c_\alpha|^2\). A two-element orbit contributes the block
\(m_\alpha\begin{pmatrix}0&1\\1&0\end{pmatrix}\), with one positive and
one negative direction. Distinct orbits are orthogonal since their
coordinates never occur together in (WP54). This proves the radical and
all three entries of (WP55).

In particular, a packet supported entirely on the critical line has
\(p=0\), and (WP54) is positive semidefinite. Its radical still contains
all higher jet directions. A distinct off-line quartet has \(f=0,p=2\),
giving (WP31). Both assertions concern the respective actual packet
algebras and the exact orders in (WP53).

\subsection{10. The exact zero-side localization and analytic
scope}\label{the-exact-zero-side-localization-and-analytic-scope}

Index the distinct nontrivial zeros of \(\xi\) by \(\rho'\), and let
\(m_{\rho'}\) be their multiplicities. For the entire functions (WP6),
form the discrete zero pairing \[
\mathscr W_0(F_P,F_Q)
=\sum_{\rho'}m_{\rho'}
\overline{F_P(\overline{\theta(\rho')})}
F_Q(\theta(\rho')).
\tag{WP56}
\] This particular sum has only four potentially nonzero summands.
Indeed, when \(\rho'\notin\mathcal R\), (WP4) gives
\(F_Q(\theta(\rho'))=q_Q(\rho')v_h(\rho')=0\). Thus no issue of
summation order remains in (WP56). At the quartet, the multiplicities
equal \(m\), and (WP8) gives the exact identity \[
\begin{aligned}
\mathscr W_0(F_P,F_Q)
&=m\sum_{\alpha\in\mathcal R}
\overline{q_P(\sigma\alpha)v_h(\sigma\alpha)}
q_Q(\alpha)v_h(\alpha)\\
&=B_h(\mathcal C[P],\mathcal C[Q]).
\end{aligned}
\tag{WP57}
\] The bracketed arguments on the second line are classes in \(E_z\),
mapped into \(E_h\) by (WP10). If a polynomial is changed by a multiple
of \(H\), its amplitude has all its first \(m\) jets zero at each
Fourier root. Therefore (WP57) descends to \(E_z\) exactly as claimed.
In fact the form already factors through the smaller value quotient
\(\mathbb C[z]/\prod_\alpha(z-\beta_\alpha)\); the explicit kernel of
the map from \(E_z\) to that quotient is \(\mathcal C^{-1}J\).

For \(m>1\), the multiplicity in (WP56) multiplies a value. It does not
differentiate \(F_P\) or \(F_Q\). This is why the correct inertia on the
full \(4m\)-dimensional packet is (WP31), and why the higher jet
directions appear as its precisely computed radical. Their maps remain
(WP21), (WP38)--(WP44), and (WP48).

The proof above does not identify (WP56) with pole, Gamma, and prime
terms for this test family. Such an identity involves its analytic test
class and convergence statements; none follows from the finite quotient
alone. The finite algebra supplies a complete target for that
calculation: the Hermitian form (WP28), with the explicit negative
representatives (WP32) or (WP36), the complete radical (WP23), and the
exact relation to the native amplitude multiplier (WP26). No assertion
here establishes positivity of the prime-side expression or a conclusion
about the existence of an off-line zero of \(\xi\).

\subsection{Exact residue-duality
continuation}\label{exact-residue-duality-continuation}

\href{COLLISION_RESIDUE_DUALITY_DERIVATION.md}{RD1--RD31b} proves the
complete Jacobian map from perfect residue duality to the retained
regular trace, identifies the actual packet radical with cotangent
cohomology, and computes the supported endpoint compensation. The
existing formulas and all their coordinates remain unchanged.

\end{document}
